\chapter{Mekanisme Reaksi Kompleks}\label{ch:b3:complex-mechanisms}

Larutkan garam kromium(III) dalam air yang dilabeli oksigen-18, dan berhari-hari
berlalu sebelum air berlabel itu menggantikan molekul-molekul air yang terikat
pada logam; dengan garam tembaga(II), pertukaran itu sudah selesai dalam
waktu pencampuran. Reaksinya sama, yaitu satu molekul air meninggalkan ion
logam dan molekul lain mengambil tempatnya, tetapi lajunya sangat berbeda
dari satu ion ke ion lain, dan jumlah elektron $d$ meramalkan sebagian besar
variasi itu. Cara ligan digantikan menentukan bagaimana obat antikanker
platina dibuat dan bagaimana obat itu bekerja; cara elektron berpindah
antara ion-ion logam menentukan kecepatan respirasi dan kecepatan setiap
baterai. Bab ini membahas dua golongan reaksi kompleks: substitusi dan
transfer elektron, yang terakhir disertai teori Rudolph Marcus.

\begin{recall}
Jilid Tahun 2 menjelaskan kompleks dan dentisitas ligan, ligan penjembatan,
isomer $fac$ dan $mer$, pertukaran ligan sebagai tahap elementer, serta energi
stabilisasi medan kristal dengan konfigurasi spin tinggi dan spin rendah;
Jilid Tahun 1 menjelaskan mekanisme beserta pendekatan keadaan tunak dan
pendekatan tahap penentu laju. \Cref{ch:b3:rate-theories} memberikan \omterm{thm:b3:rate-theories:eyring}{persamaan Eyring} dan
makna $\Delta^{\ddagger}S^\circ$, \cref{ch:b3:electrode-kinetics} transfer elektron pada
elektrode, dan \cref{ch:b3:complex-spectra-magnetism} \omterm{def:b3:complex-spectra-magnetism:ligand-field-term}{suku-suku medan ligan} dan
efek Jahn--Teller.
\end{recall}

\section{Labilitas dan keinertan}

\begin{definition}[Kompleks labil dan inert]\label{def:b3:complex-mechanisms:labile}
Sebuah kompleks disebut \emph{labil}\index{kompleks labil} jika ligan-ligannya
digantikan dengan cepat, dan disebut \emph{kompleks yang inert secara kinetika}\index{kompleks yang inert secara kinetika}
jika ligan-ligannya digantikan dengan lambat; mengikuti Taube, sebuah kompleks
disebut labil jika reaksinya dengan ligan pada \qty{0.1}{mol/L} dan \qty{25}{\celsius} selesai dalam sekitar
satu menit. Labilitas adalah sifat kinetika, yang tidak berkaitan dengan
kestabilan termodinamika kompleks itu.
\end{definition}

Pertukaran air, $\mathrm{[M(H_2O)_6]^{n+} + H_2O^{\ast} \to [M(H_2O)_5(H_2O^{\ast})]^{n+} + H_2O}$, adalah substitusi
yang paling sederhana dan menetapkan sebuah skala. Ion-ion golongan utama
bertukar lebih cepat jika ukurannya lebih besar dan muatannya lebih kecil.
Di antara ion logam transisi, ion-ion berkonfigurasi $d^3$
(\ce{Cr^{3+}}) atau $d^6$ spin rendah (\ce{Co^{3+}}, \ce{Rh^{3+}}, \ce{Ir^{3+}}) bersifat inert, demikian pula
sebagian besar logam baris kedua dan ketiga, sedangkan $d^9$ (\ce{Cu^{2+}}) dan $d^4$ spin tinggi (\ce{Cr^{2+}}),
yang terdistorsi oleh efek Jahn--Teller, sangat \omterm{def:b3:complex-mechanisms:labile}{labil}.

\begin{proposition}[Energi aktivasi medan ligan]\label{prop:b3:complex-mechanisms:lfae}
Jika sebuah substitusi melalui keadaan transisi piramida segi empat berkoordinasi
lima, hilangnya stabilisasi medan ligan dalam perjalanan ke sana (energi aktivasi
medan ligan), dalam model paling sederhana yang hanya memakai $\sigma$, paling besar untuk $d^6$ spin rendah
($4\,Dq$) dan berikutnya untuk $d^3$ dan $d^8$ ($2\,Dq$), nol untuk $d^0$, $d^5$ spin tinggi, dan $d^{10}$, serta negatif
untuk $d^4$ dan $d^9$ spin tinggi.
\end{proposition}

\begin{proof}
Modelkan setiap ligan sebagai penaik orbital $d$ sebesar $e_\sigma$ kali kuadrat amplitudo
sudut orbital itu ke arah ligan, yang dinormalisasi menjadi 1 pada sumbu cuping: ligan
aksial memberi $d_{z^2}$ 1, ligan ekuatorial memberi $d_{z^2}$ $\frac14$ dan $d_{x^2-y^2}$ $\frac34$; orbital-orbital $t_{2g}$ mengarah
di antara ligan dan tidak mendapat apa-apa. Oktahedron memberi $d_{z^2}$ dan $d_{x^2-y^2}$ masing-masing $3e_\sigma$
($\Delta_{\mathrm o} = 3e_\sigma = 10Dq$); piramida segi empat, tanpa satu ligan aksial, memberi $d_{z^2}$ $2e_\sigma$ dan
$d_{x^2-y^2}$ $3e_\sigma$. Stabilisasi $n$ elektron adalah energinya dikurangi $n/5$ kali jumlah
kelima tingkat (pusat beratnya). Untuk $d^6$ spin rendah ($t_{2g}^6$): $0 - \frac65 \times 6e_\sigma$ dalam
oktahedron, $0 - \frac65 \times 5e_\sigma$ dalam piramida: kehilangan sebesar $1.2e_\sigma = 4Dq$. Untuk $d^3$: $\frac35e_\sigma =
2Dq$; untuk $d^8$ ($t_{2g}^6e_g^2$): $(5 - \frac85 \times 5) - (6 - \frac85 \times 6) = 0.6e_\sigma$. Untuk $d^9$: $(7 - 9) - (9 -
10.8) = -0.2e_\sigma$.
\end{proof}

\begin{omfigure}
\begin{tikzpicture}
\begin{axis}[width=0.72\linewidth, height=5.4cm, xmin=-0.6, xmax=10.6, ymin=-1.2, ymax=4.6,
  axis x line=bottom, axis y line=left, ycomb, xtick={0,1,2,3,4,5,6,7,8,9,10},
  xlabel={$n$ ($d^n$)}, ylabel={energi aktivasi / $Dq$}, tick label style={font=\small},
  label style={font=\small}]
\draw[black!40] (axis cs:-0.6,0) -- (axis cs:10.6,0);
\addplot[omDef, line width=4pt, xshift=-2.5pt] table[x=n, y=hs] {figdata/out/bachelor-3-ligand-field-activation.dat};
\addplot[omThm, line width=4pt, xshift=2.5pt, unbounded coords=discard] table[x=n, y=ls] {figdata/out/bachelor-3-ligand-field-activation.dat};
\end{axis}
\end{tikzpicture}

{\small Energi aktivasi medan ligan untuk jalur disosiatif melalui piramida
segi empat, dalam model yang hanya memakai $\sigma$ pada proposisi itu (biru: spin tinggi;
merah: spin rendah). Nilai-nilai terbesar,
untuk $d^6$ spin rendah, $d^3$, dan $d^8$, sesuai dengan ion-ion yang inert atau lambat (\ce{Co^{3+}}, \ce{Cr^{3+}}, \ce{Ni^{2+}});
nilai-nilai negatif, untuk $d^4$ dan $d^9$, sesuai dengan ion-ion Jahn--Teller yang sangat \omterm{def:b3:complex-mechanisms:labile}{labil}.}
\end{omfigure}

\begin{definition}[Volume aktivasi]\label{def:b3:complex-mechanisms:activation-volume}
Yang disebut \emph{volume aktivasi}\index{volume aktivasi} $\Delta V^{\ddagger}$ sebuah tahap
elementer adalah selisih antara volume molar parsial \omterm{def:b3:rate-theories:activated-complex}{kompleks teraktivasi} dan
volume molar parsial pereaksi-pereaksinya.
\end{definition}

\begin{proposition}[Ketergantungan tetapan laju pada tekanan]\label{prop:b3:complex-mechanisms:pressure}
Pada suhu tetap, $\displaystyle\Bigl(\frac{\partial\ln k}{\partial p}\Bigr)_T = -\frac{\Delta V^{\ddagger}}{RT}$.
\end{proposition}

\begin{proof}
Menurut \omterm{thm:b3:rate-theories:eyring}{persamaan Eyring} $\ln k = \text{tetapan} - \Delta^{\ddagger}G/RT$, dan $(\partial\Delta^{\ddagger}G/\partial p)_T = \Delta V^{\ddagger}$, seperti untuk
setiap energi Gibbs.
\end{proof}

Tahap yang melepaskan sebuah ligan memiliki $\Delta V^{\ddagger} > 0$; tahap yang memasukkan sebuah ligan memiliki $\Delta V^{\ddagger}
< 0$. Jika diukur pada beberapa ratus megapascal, \omterm{def:b3:complex-mechanisms:activation-volume}{volume aktivasi} merupakan
diagnosis mekanisme substitusi yang paling langsung, yang lebih sedikit
dikaburkan oleh solvasi daripada \omterm{def:b3:rate-theories:activation}{entropi aktivasi}.

\section{Mekanisme substitusi}

\begin{definition}[Mekanisme disosiatif, asosiatif, dan pertukaran]\label{def:b3:complex-mechanisms:mechanisms}
Dalam \emph{mekanisme disosiatif}\index{mekanisme disosiatif} (D), ligan yang
pergi lepas lebih dahulu, sehingga terbentuk zat antara dengan bilangan
koordinasi yang lebih rendah; dalam \emph{mekanisme asosiatif}\index{mekanisme asosiatif} (A), ligan
yang masuk terikat lebih dahulu, sehingga terbentuk zat antara dengan bilangan
koordinasi yang lebih tinggi; dalam \emph{mekanisme pertukaran}\index{mekanisme pertukaran} (I),
kedua ligan bertukar dalam satu tahap, dengan karakter pemutusan ikatan ($I_d$)
atau pembentukan ikatan ($I_a$).
\end{definition}

\begin{proposition}[Hukum laju mekanisme D dan A]\label{prop:b3:complex-mechanisms:d-a-laws}
Untuk $\mathrm{ML_5X + Y \to ML_5Y + X}$: mekanisme D ($\mathrm{ML_5X} \rightleftharpoons \mathrm{ML_5} + \mathrm X$, $k_1$, $k_{-1}$; $\mathrm{ML_5} + \mathrm Y
\to \mathrm{ML_5Y}$, $k_2$) memberikan
\[
  v = \frac{k_1k_2[\mathrm{ML_5X}][\mathrm Y]}{k_{-1}[\mathrm X] + k_2[\mathrm Y]},
\]
yang berorde satu terhadap kompleks dan tidak bergantung pada $[\mathrm Y]$ pada $[\mathrm Y]$ yang tinggi; mekanisme A
melalui zat antara berkoordinasi tujuh dalam keadaan tunak memberikan $v =
k[\mathrm{ML_5X}][\mathrm Y]$, yang berorde satu terhadap masing-masing.
\end{proposition}

\begin{proof}
D: keadaan tunak untuk $\mathrm{ML_5}$, $k_1[\mathrm{ML_5X}] = (k_{-1}[\mathrm X] + k_2[\mathrm Y])[\mathrm{ML_5}]$, dan $v = k_2[\mathrm{ML_5}][\mathrm Y]$.
Pada $[\mathrm Y]$ yang tinggi, $v \to k_1[\mathrm{ML_5X}]$. A: $\mathrm{ML_5X} + \mathrm Y \rightleftharpoons \mathrm{ML_5XY}$ ($k_a$, $k_{-a}$), $\mathrm{ML_5XY} \to
\mathrm{ML_5Y} + \mathrm X$ ($k_b$); keadaan tunak memberikan $v = \frac{k_ak_b}{k_{-a} + k_b}[\mathrm{ML_5X}][\mathrm Y]$.
\end{proof}

Di dalam air, ligan yang masuk biasanya hadir pada konsentrasi yang jauh lebih
rendah daripada pelarutnya, dan tahap pertama yang sebenarnya adalah
pertemuan kompleks dengan ligan itu.

\begin{definition}[Mekanisme Eigen--Wilkins]\label{def:b3:complex-mechanisms:eigen-wilkins}
Yang disebut \emph{mekanisme Eigen--Wilkins}\index{mekanisme Eigen--Wilkins}
substitusi sebuah kompleks akua adalah prakesetimbangan cepat yang membentuk
\emph{kompleks sfer luar}\index{kompleks sfer luar}, tempat ligan yang masuk
berada di samping kompleks, di luar sfer koordinasinya, yang diikuti oleh
pertukaran sebuah molekul air dengan ligan itu sebagai tahap penentu laju.
\end{definition}

\begin{theorem}[Hukum laju Eigen--Wilkins]\label{thm:b3:complex-mechanisms:eigen-wilkins}
Dengan tetapan kesetimbangan sfer luar $K_{\mathrm{os}}$ dan tetapan laju pertukaran
$k_i$, tetapan laju orde satu semu yang teramati pada $[\mathrm Y]$ berlebih adalah
\[
  k_{\mathrm{obs}} = \frac{K_{\mathrm{os}}k_i[\mathrm Y]}{1 + K_{\mathrm{os}}[\mathrm Y]} .
\]
\end{theorem}

\begin{proof}
Kompleks itu terdistribusi antara bentuk bebas dan bentuk sfer luar, $[\mathrm{M\cdot Y}] = K_{\mathrm{os}}[\mathrm M][\mathrm Y]$,
dengan total $[\mathrm M]_{\mathrm{tot}} = [\mathrm M](1 + K_{\mathrm{os}}[\mathrm Y])$; lajunya adalah $k_i[\mathrm{M\cdot Y}] = k_iK_{\mathrm{os}}[\mathrm Y][\mathrm M]_{\mathrm{tot}}/(1 + K_{\mathrm{os}}[\mathrm Y])$.
\end{proof}

Pada $[\mathrm Y]$ yang rendah, reaksinya berorde dua dengan $k = K_{\mathrm{os}}k_i$; $K_{\mathrm{os}}$, yang diperkirakan dari
muatan dan jarak pendekatan terdekat, hampir sama untuk semua ligan yang
muatannya sama, sehingga $k_i$, yaitu tahap pertukaran air, mengendalikan
laju: substitusi pada ion akua hampir tidak bergantung pada ligan yang
masuk, dan inilah tanda pertukaran disosiatif.

\begin{definition}[Mekanisme basa konjugat]\label{def:b3:complex-mechanisms:conjugate-base}
Yang disebut \emph{mekanisme basa konjugat}\index{mekanisme basa konjugat} hidrolisis
basa sebuah kompleks amin, $\mathrm{[Co(NH_3)_5X]^{2+} + OH^- \to [Co(NH_3)_5(OH)]^{2+} + X^-}$, adalah
deprotonasi sebuah ligan amin dalam prakesetimbangan yang cepat, yang diikuti oleh
pelepasan $\mathrm X^-$ secara disosiatif dari basa konjugat amido, yang ligan $\mathrm{NH_2^-}$-nya
melabilkan kompleks itu.
\end{definition}

\begin{proposition}[Hukum laju hidrolisis basa]\label{prop:b3:complex-mechanisms:conjugate-base}
Dengan tetapan deprotonasi $K$ dan tetapan laju disosiasi $k$ basa
konjugatnya, $v = \dfrac{kK[\text{kompleks}]_{\mathrm{tot}}[\ce{OH-}]}{1 + K[\ce{OH-}]}$, yang berorde dua pada $[\ce{OH-}]$ yang rendah.
\end{proposition}

\begin{proof}
Seperti pada hukum Eigen--Wilkins: kompleks itu terdistribusi antara bentuk
asam dan bentuk basa konjugatnya dengan perbandingan $1 : K[\ce{OH-}]$, dan hanya bentuk yang kedua yang bereaksi.
\end{proof}

Hukum laju saja tidak membedakan mekanisme ini dari serangan langsung oleh
\ce{OH-}; buktinya adalah bahwa hidrolisis basa memerlukan sebuah proton N--H (kompleks yang
tidak memilikinya tidak memperlihatkan hidrolisis itu) dan bahwa proton-proton
amin bertukar dengan pelarut jauh lebih cepat daripada hidrolisisnya.

\begin{method}[Mendiagnosis mekanisme substitusi]\label{met:b3:complex-mechanisms:diagnose}
\begin{enumerate}
  \item Ketergantungan pada gugus yang masuk: laju yang tidak bergantung pada Y (atau yang jenuh),
        disosiatif; laju yang sebanding dengan $[\mathrm Y]$ dan peka terhadap jenis Y,
        asosiatif.
  \item \omterm{def:b3:complex-mechanisms:activation-volume}{Volume aktivasi}: positif, disosiatif; negatif, asosiatif.
  \item \omterm{def:b3:rate-theories:activation}{Entropi aktivasi}: positif untuk D, negatif untuk A (dengan hati-hati, karena
        solvasi ikut berkontribusi).
  \item Ketergantungan pada gugus yang pergi: laju yang mengikuti kekuatan ikatan M--X
        menunjukkan pemutusan ikatan dalam keadaan transisi.
\end{enumerate}
\end{method}

\begin{omfigure}
\begin{tikzpicture}
\begin{axis}[name=pd, omprofile, width=0.46\linewidth, height=4.6cm, xmin=0, xmax=10, ymin=0, ymax=10,
  xlabel={koordinat reaksi}, ylabel={energi}, title={\small D: $\mathrm{ML_5X} \to \mathrm{ML_5} + \mathrm X \to \mathrm{ML_5Y}$}]
\addplot[omDef, very thick, smooth] coordinates {(0,2) (1.5,2.1) (3,7.5) (4.2,5.6) (5.0,5.4) (5.8,5.6) (7,7.0) (8.5,1.6) (10,1.5)};
\node[font=\scriptsize, anchor=north, align=center] at (axis cs:5,5.2) {ML$_5$ (BK 5)};
\end{axis}
\begin{axis}[at={(pd.east)}, anchor=west, xshift=1.0cm, omprofile, width=0.46\linewidth, height=4.6cm,
  xmin=0, xmax=10, ymin=0, ymax=10, xlabel={koordinat reaksi}, ylabel={energi},
  title={\small A: $\mathrm{ML_5X} + \mathrm Y \to \mathrm{ML_5XY} \to \mathrm{ML_5Y}$}]
\addplot[omThm, very thick, smooth] coordinates {(0,2) (1.5,2.1) (3,6.5) (4.2,4.4) (5.0,4.2) (5.8,4.4) (7,7.2) (8.5,1.6) (10,1.5)};
\node[font=\scriptsize, anchor=north, align=center] at (axis cs:5,4.0) {ML$_5$XY (BK 7)};
\end{axis}
\end{tikzpicture}

{\small Profil energi substitusi disosiatif dan asosiatif
(skematis): keduanya melalui sebuah zat antara, yang bilangan koordinasinya
(BK) lebih rendah pada D dan lebih tinggi pada A. Pada \omterm{def:b3:complex-mechanisms:mechanisms}{mekanisme pertukaran},
sumurnya lenyap dan hanya tersisa satu keadaan transisi.}
\end{omfigure}

Kompleks oktahedral yang bersubstitusi melalui piramida segi empat umumnya
mempertahankan konfigurasinya ($cis$ tetap $cis$); kompleks yang melalui
bipiramida trigonal, seperti banyak kompleks kobalt(III) pada hidrolisis basa, dapat memberikan
campuran produk $cis$ dan $trans$.

\section{Substitusi segi empat datar dan efek trans}

Kompleks segi empat datar ion-ion $d^8$ (\ce{Pt^{2+}}, \ce{Pd^{2+}}, \ce{Au^{3+}}, \ce{Ni^{2+}} dengan ligan kuat)
tidak jenuh secara koordinasi: ligan yang masuk dapat mendekat sepanjang arah
aksial yang kosong. Substitusinya bersifat asosiatif, melalui keadaan transisi
bipiramida trigonal berkoordinasi lima.

\begin{proposition}[Hukum laju dua suku]\label{prop:b3:complex-mechanisms:square-planar}
Substitusi $\mathrm{[PtL_3X] + Y \to [PtL_3Y] + X}$ dalam pelarut S yang berkoordinasi mengikuti
$k_{\mathrm{obs}} = k_1 + k_2[\mathrm Y]$ pada Y berlebih.
\end{proposition}

\begin{proof}
Dua jalur asosiatif yang paralel: serangan langsung oleh Y ($k_2[\mathrm Y]$), dan serangan oleh
pelarut, yang hadir dalam jumlah sangat berlebih (orde satu semu $k_1$), yang memberikan $[\mathrm{PtL_3S}]$, yang
kemudian dengan cepat diubah oleh Y. Jalur-jalur orde satu yang paralel saling menjumlah.
\end{proof}

\begin{definition}[Efek trans, pengaruh trans]\label{def:b3:complex-mechanisms:trans-effect}
Yang disebut \emph{efek trans}\index{efek trans} adalah percepatan substitusi
sebuah ligan oleh ligan yang berada pada posisi trans terhadapnya: sebuah efek kinetika, pada keadaan transisi.
Yang disebut \emph{pengaruh trans}\index{pengaruh trans} adalah pelemahan, dalam keadaan
dasar, ikatan yang trans terhadap sebuah ligan, yang terlihat pada panjang ikatan dan frekuensi
vibrasi: sebuah efek termodinamika.
\end{definition}

\omterm{def:b3:complex-mechanisms:trans-effect}{Efek trans} mengikuti urutan $\ce{CO}$, $\ce{CN-}$, $\ce{C2H4} > \ce{PR3}$, $\ce{H-} > \ce{CH3-} > \ce{I-}$, $\ce{SCN-} >
\ce{Br-} > \ce{Cl-} > \ce{NH3}$, $\ce{H2O}$. Donor $\sigma$ yang kuat melemahkan ikatan trans (\omterm{def:b3:complex-mechanisms:trans-effect}{pengaruh trans});
akseptor $\pi$ yang kuat menstabilkan keadaan transisi berkoordinasi lima dengan menarik
\omterm{def:b3:computational-chemistry:dft}{kerapatan elektron} dari logam. Keduanya menaikkan laju.

\begin{method}[Merencanakan sintesis isomer platina(II)]\label{met:b3:complex-mechanisms:pt-synthesis}
\begin{enumerate}
  \item Pada setiap tahap, ligan yang digantikan adalah ligan yang trans terhadap ligan
        dengan \omterm{def:b3:complex-mechanisms:trans-effect}{efek trans} tertinggi yang ada.
  \item Urutkan penambahan-penambahan sedemikian rupa sehingga aturan ini menghasilkan isomer yang dikehendaki.
\end{enumerate}
\end{method}

\begin{omfigure}
\setchemfig{atom sep=2.2em}
\schemestart
\chemfig{Cl-Pt(-[2]Cl)(-[6]Cl)-Cl}
\arrow{->[\ce{NH3}]}
\chemfig{Cl-Pt(-[2]Cl)(-[6]Cl)-NH_3}
\arrow{->[\ce{NH3}]}
\chemfig{Cl-Pt(-[2]NH_3)(-[6]Cl)-NH_3}
\schemestop

\bigskip
\schemestart
\chemfig{H_3N-Pt(-[2]NH_3)(-[6]NH_3)-NH_3}
\arrow{->[\ce{Cl-}]}
\chemfig{H_3N-Pt(-[2]NH_3)(-[6]NH_3)-Cl}
\arrow{->[\ce{Cl-}]}
\chemfig{Cl-Pt(-[2]NH_3)(-[6]NH_3)-Cl}
\schemestop

{\small \omterm{def:b3:complex-mechanisms:trans-effect}{Efek trans} dalam praktik (muatan tidak dituliskan). Atas: dari $\ce{[PtCl4]^{2-}}$, amonia
kedua menggantikan klorida yang trans terhadap klorida (Cl$^-$ memiliki \omterm{def:b3:complex-mechanisms:trans-effect}{efek trans} yang lebih besar),
sehingga terbentuk isomer $cis$, yaitu cisplatin. Bawah: dari $\ce{[Pt(NH3)4]^{2+}}$, klorida kedua
menggantikan amonia yang trans terhadap klorida pertama, sehingga terbentuk isomer $trans$.}
\end{omfigure}

\section{Transfer elektron}

\begin{definition}[Transfer elektron sfer luar dan sfer dalam]\label{def:b3:complex-mechanisms:electron-transfer}
Pada \emph{transfer elektron sfer luar}\index{transfer elektron sfer luar},
elektron berpindah antara dua kompleks yang sfer koordinasinya tetap utuh;
pada \emph{transfer elektron sfer dalam}\index{transfer elektron sfer dalam},
elektron berpindah melalui ligan yang menjembatani kedua logam dalam sebuah
kompleks dwiinti yang sementara. Yang disebut \emph{reaksi swapertukaran}\index{reaksi swapertukaran} adalah
transfer elektron antara dua bilangan oksidasi pasangan yang sama, tanpa
perubahan kimia neto, misalnya $\ce{[Fe(H2O)6]^{3+} + [Fe(H2O)6]^{2+}}$.
\end{definition}

Henry Taube memperlihatkan jalur sfer dalam pada 1953 dengan reaksi yang dipilih
sedemikian rupa sehingga produknya menceritakan kisahnya. Kobalt(III) inert, kobalt(II) \omterm{def:b3:complex-mechanisms:labile}{labil}; kromium(II)
\omterm{def:b3:complex-mechanisms:labile}{labil}, kromium(III) inert. Ketika $\ce{[Co(NH3)5Cl]^{2+}}$ direduksi oleh $\ce{[Cr(H2O)6]^{2+}}$ dalam
asam, seluruh kromium(III) yang terbentuk membawa klorida, sebagai $\ce{[Cr(H2O)5Cl]^{2+}}$, bahkan dengan
hadirnya klorida bebas: klorida itu pasti telah menjembatani kedua logam,
terikat pada kromium(III) yang inert begitu elektron berpindah, dan dilepaskan
oleh kobalt(II) yang \omterm{def:b3:complex-mechanisms:labile}{labil}.

\begin{omfigure}
\begin{tikzpicture}[font=\small]
  \begin{scope}[scale=0.8]
    \pic[transform shape] (a) at (0,0) {omoctahedron};
    \pic[transform shape] (b) at (2.0,0) {omoctahedron};
  \end{scope}
  \node[anchor=north east] at (0.3,-1.0) {$\mathrm{Co^{III}(NH_3)_5}$};
  \node[anchor=north west] at (1.3,-1.0) {$\mathrm{Cr^{II}(H_2O)_5}$};
  \node[omThm, anchor=south] at (0.8,0.15) {Cl};
  \draw[->, thick, omDef] (1.6,1.75) to[bend right=40] node[midway, above] {e$^-$} (0,1.75);
\end{tikzpicture}

{\small Zat antara berjembatan Taube (skematis): klorida, yang masih terikat pada
kobalt(III), memasuki sfer koordinasi kromium(II) dengan menggantikan sebuah molekul
air; elektron menyeberang melalui jembatan itu, dan klorida tetap berada pada
kromium(III) yang kini inert.}
\end{omfigure}

Sebuah elektron melompat dalam waktu sekitar $10^{-16}$ s, jauh lebih cepat daripada gerak inti: menurut
\omterm{thm:b3:electronic-spectroscopy:franck-condon}{prinsip Franck--Condon} (\cref{ch:b3:electronic-spectroscopy}), transfer itu
hanya dapat terjadi pada konfigurasi inti yang membuat pereaksi dan produk memiliki
energi yang sama. Pada swapertukaran besi, ikatan \ce{Fe^{III}}--O lebih pendek daripada
ikatan \ce{Fe^{II}}--O; sebelum elektron dapat berpindah, kedua kompleks harus terdistorsi ke
geometri antara yang sama, dengan mengorbankan energi.

\section{Teori Marcus}

\begin{definition}[Energi reorganisasi]\label{def:b3:complex-mechanisms:reorganisation}
Yang disebut \emph{energi reorganisasi}\index{energi reorganisasi} $\lambda$ sebuah
transfer elektron adalah energi yang diperlukan untuk membawa pereaksi, tanpa
memindahkan elektronnya, ke konfigurasi inti (panjang ikatan kompleks dan
orientasi pelarut di sekelilingnya) milik produk. Energi ini merupakan jumlah
bagian dalam, dari ikatan, dan bagian luar, dari pelarut.
\end{definition}

\begin{theorem}[Persamaan Marcus]\label{thm:b3:complex-mechanisms:marcus}
Jika energi Gibbs pereaksi dan produk berupa parabola yang kelengkungannya sama
dalam sebuah koordinat reaksi $x$ (0 pada kesetimbangan pereaksi, 1 pada kesetimbangan
produk), \omterm{def:b3:rate-theories:activation}{energi Gibbs aktivasi} sebuah transfer elektron dengan energi Gibbs
reaksi standar $\Delta G^\circ$ dan \omterm{def:b3:complex-mechanisms:reorganisation}{energi reorganisasi} $\lambda$ adalah
\[
  \Delta G^{\ddagger} = \frac{(\lambda + \Delta G^\circ)^2}{4\lambda} .
\]
Inilah \emph{persamaan Marcus}\index{persamaan Marcus}.
\end{theorem}

\begin{proof}
Tuliskan $G_R = \lambda x^2$ dan $G_P = \lambda(x - 1)^2 + \Delta G^\circ$: kelengkungannya ditetapkan oleh $G_R(1) = \lambda$, yaitu
definisi $\lambda$. Transfer terjadi di tempat kedua kurva bersilangan: $\lambda x^2 = \lambda x^2 - 2\lambda x + \lambda +
\Delta G^\circ$, sehingga $x^\ast = (\lambda + \Delta G^\circ)/2\lambda$, dan $\Delta G^{\ddagger} = G_R(x^\ast) = (\lambda + \Delta G^\circ)^2/4\lambda$.
\end{proof}

\begin{corollary}[Daerah terbalik]\label{cor:b3:complex-mechanisms:inverted}
Pada $\lambda$ tetap, laju bertambah ketika reaksi menjadi makin eksergonik sampai
$-\Delta G^\circ = \lambda$, tempat $\Delta G^{\ddagger} = 0$; di luar titik itu, laju berkurang lagi.
\end{corollary}

\begin{proof}
$\Delta G^{\ddagger}$ adalah parabola dalam $\Delta G^\circ$ dengan minimum, nol, pada $\Delta G^\circ = -\lambda$; untuk $\Delta G^\circ < -\lambda$ nilainya tumbuh
lagi.
\end{proof}

\begin{definition}[Daerah terbalik]\label{def:b3:complex-mechanisms:inverted}
Yang disebut \emph{daerah terbalik}\index{daerah terbalik} transfer elektron adalah
rentang gaya dorong $-\Delta G^\circ > \lambda$ tempat laju berkurang ketika reaksi menjadi
makin eksergonik.
\end{definition}

\begin{omfigure}
\begin{tikzpicture}
\begin{axis}[name=mp, width=0.47\linewidth, height=5.6cm, xmin=-0.6, xmax=2.2, ymin=-160, ymax=120,
  axis lines=left, xlabel={koordinat reaksi $x$}, ylabel={$G$ / \unit{kJ.mol^{-1}}},
  tick label style={font=\small}, label style={font=\small}, xtick={0,1,2}]
\addplot[black, thick] table[x=x, y=GR] {figdata/out/bachelor-3-marcus-parabolas.dat};
\addplot[omDef!50, thick] table[x=x, y=P0] {figdata/out/bachelor-3-marcus-parabolas.dat};
\addplot[omDef, thick] table[x=x, y=P50] {figdata/out/bachelor-3-marcus-parabolas.dat};
\addplot[omThm, thick] table[x=x, y=P100] {figdata/out/bachelor-3-marcus-parabolas.dat};
\addplot[omThm!60, thick, dashed] table[x=x, y=P150] {figdata/out/bachelor-3-marcus-parabolas.dat};
\node[font=\scriptsize, anchor=north] at (axis cs:0,-2) {R};
\end{axis}
\begin{axis}[at={(mp.east)}, anchor=west, xshift=1.4cm, width=0.47\linewidth, height=5.6cm,
  xmin=0, xmax=250, ymin=0, ymax=11.5, axis lines=left, xlabel={$-\Delta G^\circ$ / \unit{kJ.mol^{-1}}},
  ylabel={$\log(k/\unit{s^{-1}})$}, tick label style={font=\small}, label style={font=\small}]
\addplot[omDef, very thick] table[x=mdG, y=logk] {figdata/out/bachelor-3-marcus-rate.dat};
\draw[black!40, dashed] (axis cs:100,0) -- (axis cs:100,11);
\node[font=\scriptsize, anchor=south east] at (axis cs:95,1) {normal};
\node[font=\scriptsize, anchor=south west] at (axis cs:105,1) {terbalik};
\end{axis}
\end{tikzpicture}

{\small Teori Marcus dengan $\lambda = \qty{100}{kJ/mol}$ (model). Kiri: parabola pereaksi
(hitam) dan parabola produk untuk $\Delta G^\circ = 0$ dan $-50$ (biru), $-100 = -\lambda$ (merah), dan
\qty{-150}{kJ/mol} (putus-putus, \omterm{def:b3:complex-mechanisms:inverted}{daerah terbalik}); titik
silangnya, yaitu keadaan transisi, turun sampai dasar kurva pereaksi pada
$-\Delta G^\circ = \lambda$ dan naik lagi di luar titik itu. Kanan: logaritma tetapan laju
(faktor praeksponensial model \qty{e11}{s^{-1}}) melewati maksimum pada $-\Delta G^\circ = \lambda$.}
\end{omfigure}

\begin{proposition}[Energi reorganisasi sfer luar]\label{prop:b3:complex-mechanisms:outer-reorganisation}
Untuk dua pereaksi bulat berjari-jari $a_1$, $a_2$ pada jarak $d$ dalam pelarut dengan
indeks bias $n$ dan permitivitas relatif $\varepsilon_s$, bagian pelarut dari $\lambda$
sebanding dengan $\bigl(\frac{1}{2a_1} + \frac{1}{2a_2} - \frac1d\bigr)\bigl(\frac{1}{n^2} - \frac{1}{\varepsilon_s}\bigr)$.
\end{proposition}

\admitted

Air, yang sangat polar ($\varepsilon_s$ besar) tetapi indeks biasnya biasa saja, memberikan
reorganisasi pelarut yang besar; pelarut nonpolar memberikan reorganisasi yang kecil. Pereaksi yang
besar lebih sedikit bereorganisasi: itulah sebabnya protein transfer elektron
membenamkan situs logamnya di dalam protein.

\begin{theorem}[Hubungan silang Marcus]\label{thm:b3:complex-mechanisms:cross-relation}
Untuk reaksi silang $\mathrm{A_{ox} + B_{red} \to A_{red} + B_{ox}}$ dengan tetapan kesetimbangan $K_{12}$, antara pasangan-pasangan
yang tetapan laju swapertukarannya $k_{11}$ dan $k_{22}$,
\[
  k_{12} = \sqrt{k_{11}k_{22}K_{12}f}, \qquad \ln f = \frac{(\ln K_{12})^2}{4\ln(k_{11}k_{22}/Z^2)},
\]
dengan $Z$ faktor frekuensi tumbukan; $f \approx 1$ jika $K_{12}$ tidak terlalu besar. Inilah
\emph{hubungan silang Marcus}\index{hubungan silang Marcus}.
\end{theorem}

\begin{proof}[Bukti parsial]
Anggaplah \omterm{def:b3:complex-mechanisms:reorganisation}{energi reorganisasi} reaksi silang adalah rata-rata \omterm{def:b3:complex-mechanisms:reorganisation}{energi reorganisasi}
kedua swapertukaran, $\lambda_{12} = (\lambda_{11} + \lambda_{22})/2$ (setiap pereaksi menyumbang setengah \omterm{def:b3:complex-mechanisms:reorganisation}{energi reorganisasi}
swapertukarannya sendiri), dan semua tetapan laju $k = Z\eu^{-\Delta G^{\ddagger}/RT}$. Maka
$\Delta G_{12}^{\ddagger} = \frac{\lambda_{12}}{4}\bigl(1 + \frac{\Delta G^\circ_{12}}{\lambda_{12}}\bigr)^2 = \frac{\lambda_{12}}{4} + \frac{\Delta G^\circ_{12}}{2} + \frac{(\Delta G^\circ_{12})^2}{4\lambda_{12}}$. Suku pertama memberikan
$\sqrt{k_{11}k_{22}}$ (karena $\Delta G^{\ddagger}_{ii} = \lambda_{ii}/4$), suku kedua $\sqrt{K_{12}}$ (dengan $\Delta G^\circ_{12} = -RT\ln K_{12}$), dan suku ketiga
faktor $f$.
\end{proof}

\begin{method}[Laju reaksi silang dari data swapertukaran]\label{met:b3:complex-mechanisms:marcus-estimate}
\begin{enumerate}
  \item Carilah tetapan swapertukaran $k_{11}$, $k_{22}$ kedua pasangan itu.
  \item Hitunglah $K_{12}$ dari potensial standar: $\ln K_{12} = nF(E^\circ_{\text{oksidator}} - E^\circ_{\text{reduktor}})/RT$.
  \item $k_{12} = \sqrt{k_{11}k_{22}K_{12}}$ (dengan $f$ jika $K_{12}$ sangat besar). Kecocokan dalam
        faktor beberapa kali dengan nilai terukur menunjukkan mekanisme sfer luar.
\end{enumerate}
\end{method}

Marcus menerbitkan teori itu pada 1956; \omterm{def:b3:complex-mechanisms:inverted}{daerah terbalik}, yang mula-mula diragukan,
teramati pada 1984 oleh Gerhard Closs dan John Miller dalam molekul-molekul yang
donor dan akseptornya dipertahankan pada jarak tetap oleh pengatur jarak yang
kaku, sehingga difusi tidak dapat menyembunyikannya.

\begin{inthelab}[Pertukaran air dengan NMR oksigen-17]
Resonansi oksigen-17 air yang terikat pada ion paramagnetik dan resonansi air
ruah sama-sama melebar karena pertukaran di antara keduanya. Dengan merekam
lebar garis sinyal air ruah terhadap suhu, dalam larutan perklorat logam itu
yang diperkaya \ce{^{17}O}, diperoleh tetapan laju pertukarannya dan, dengan probe
tekanan tinggi, \omterm{def:b3:complex-mechanisms:activation-volume}{volume aktivasinya}.
\end{inthelab}

\begin{safety}{\ghs[0.8cm]{GHS05}\ghs[0.8cm]{GHS06}\ghs[0.8cm]{GHS08}}
% ledger: ghs4:K2PtCl4
Kalium tetrakloroplatinat(II): beracun jika tertelan, menyebabkan kerusakan
mata yang serius, pemeka saluran pernapasan dan kulit. Kompleks platina
ditangani dengan sarung tangan di dalam lemari asam; cisplatin dan
analog-analognya adalah obat sitotoksik dan ditangani sebagai obat sitotoksik.
\end{safety}

\begin{history}[Taube dan Marcus]
Henry Taube menggolongkan kompleks menjadi \omterm{def:b3:complex-mechanisms:labile}{labil} atau inert pada 1952, dan
dengan percobaan kromium--kobaltnya pada 1953 membuktikan mekanisme sfer
dalam; ia menerima Hadiah Nobel Kimia 1983. Sejak 1956 Rudolph Marcus
menurunkan teori transfer elektron yang menyandang namanya, dan ia menerima
Hadiah Nobel Kimia 1992.
\end{history}

\section{Latihan}

\begin{exercise}[$\star$]\label{exo:b3:complex-mechanisms:1}
Golongkan sebagai \omterm{def:b3:complex-mechanisms:labile}{labil} atau inert, disertai alasan: $\ce{[Cr(H2O)6]^{3+}}$, $\ce{[Co(NH3)6]^{3+}}$, $\ce{[Cr(H2O)6]^{2+}}$,
$\ce{[Cu(H2O)6]^{2+}}$, $\ce{[Zn(H2O)6]^{2+}}$.
\end{exercise}

\begin{exercise}[$\star$]\label{exo:b3:complex-mechanisms:2}
Tunjukkan bahwa hukum laju D tereduksi menjadi $v = k_1[\mathrm{ML_5X}]$ pada $[\mathrm Y]$ yang tinggi dan menjadi $v = (k_1k_2/k_{-1})
[\mathrm{ML_5X}][\mathrm Y]/[\mathrm X]$ jika $k_{-1}[\mathrm X] \gg k_2[\mathrm Y]$.
\end{exercise}

\begin{exercise}[$\star$]\label{exo:b3:complex-mechanisms:3}
Ramalkan tanda $\Delta V^{\ddagger}$ dan $\Delta^{\ddagger}S^\circ$ untuk substitusi D dan substitusi A.
\end{exercise}

\begin{exercise}[$\star$]\label{exo:b3:complex-mechanisms:4}
Berikan produk $\ce{[PtCl4]^{2-}}$ dengan dua \ce{NH3}, dan produk $\ce{[Pt(NH3)4]^{2+}}$ dengan dua \ce{Cl-}.
\end{exercise}

\begin{exercise}[$\star\star$]\label{exo:b3:complex-mechanisms:5}
Untuk sebuah substitusi pada ion akua, $K_{\mathrm{os}} = \qty{2.0}{L.mol^{-1}}$ dan $k_i = \qty{3.0e4}{s^{-1}}$ (data
latihan). Hitunglah $k_{\mathrm{obs}}$ pada $[\mathrm Y] = 0.10$ dan \qty{1.0}{mol/L}, serta limitnya.
\end{exercise}

\begin{exercise}[$\star\star$]\label{exo:b3:complex-mechanisms:6}
Sebuah tetapan laju menjadi dua kali lipat ketika tekanan naik dari 0.1 menjadi \qty{100}{MPa} pada \qty{298}{K}.
Hitunglah $\Delta V^{\ddagger}$ dan tariklah kesimpulan.
\end{exercise}

\begin{exercise}[$\star\star$]\label{exo:b3:complex-mechanisms:7}
Dalam $\ce{[PtCl3(PEt3)]-}$, ikatan Pt--Cl yang trans terhadap \ce{PEt3} panjangnya \qty{2.42}{\angstrom} dan dua ikatan lainnya \qty{2.32}{\angstrom}
(data latihan). Tafsirkan, dan ramalkan klorida mana yang digantikan lebih dahulu.
\end{exercise}

\begin{exercise}[$\star\star$]\label{exo:b3:complex-mechanisms:8}
Sebutkan bukti-bukti eksperimen yang akan menunjukkan bahwa sebuah transfer elektron
berlangsung melalui sfer dalam.
\end{exercise}

\begin{exercise}[$\star\star$]\label{exo:b3:complex-mechanisms:9}
Untuk $\lambda = \qty{80}{kJ/mol}$, hitunglah $\Delta G^{\ddagger}$ untuk $\Delta G^\circ = 0$, $-20$, $-80$, dan \qty{-120}{kJ/mol}.
\end{exercise}

\begin{exercise}[$\star\star\star$]\label{exo:b3:complex-mechanisms:10}
Dua pasangan memiliki tetapan swapertukaran $k_{11} = 4.0$ dan $k_{22} = \qty{2.0e5}{L.mol^{-1}.s^{-1}}$, dan
reaksi silangnya memiliki $K_{12} = \num{1.0e4}$ (data latihan). Perkirakan $k_{12}$ ($f = 1$).
\end{exercise}

\begin{exercise}[$\star\star\star$]\label{exo:b3:complex-mechanisms:11}
Dengan model yang hanya memakai $\sigma$, hitunglah energi aktivasi medan ligan ion-ion
$d^5$, $d^7$, dan $d^8$ spin tinggi, lalu urutkan \ce{Mn^{2+}}, \ce{Co^{2+}}, dan \ce{Ni^{2+}} berdasarkan laju
pertukaran air yang diharapkan.
\end{exercise}

\begin{exercise}[$\star\star\star$]\label{exo:b3:complex-mechanisms:12}
Mengapa \omterm{def:b3:complex-mechanisms:inverted}{daerah terbalik} sulit diamati untuk reaksi antara ion-ion yang
berdifusi bersama dalam larutan, dan bagaimana molekul donor--akseptor yang kaku
mengungkapkannya?
\end{exercise}

\section{Soal: Membuat Cisplatin, Bukan Transplatin}

\begin{problem}[{Soal akhir pekan --- membuat cisplatin, bukan transplatin: ion
platina(II) segi empat datar, sintesis yang direncanakan dengan efek trans,
akuasi obat itu, dan mengapa obat itu teraktivasi di dalam sel}]
\label{pb:b3:complex-mechanisms:1}
% ledger: ghs4:K2PtCl4
Cisplatin, $cis$-$\ce{[PtCl2(NH3)2]}$, dibuat dari $\ce{K2[PtCl4]}$. Data soal: pada
\qty{37}{\celsius}, akuasi pertama, $\ce{[PtCl2(NH3)2] + H2O -> [PtCl(H2O)(NH3)2]+ + Cl-}$, memiliki
tetapan laju $k = \qty{9.0e-5}{s^{-1}}$ dan tetapan kesetimbangan $K = \qty{3.0e-3}{mol/L}$; konsentrasi
klorida sekitar \qty{100}{mM} dalam plasma darah dan \qty{4}{mM} di dalam sel.

\medskip\noindent\textbf{Bagian I --- Platina(II).}
\begin{enumerate}[beginpenalty=10000]
  \item Berikan jumlah elektron $d$ pada \ce{Pt^{2+}}.
  \item Mengapa kompleks-kompleksnya segi empat datar dan diamagnetik?
  \item Mengapa substitusinya bersifat asosiatif?
  \item Tuliskan hukum laju dua suku dan jelaskan kedua sukunya.
  \item Apakah obat itu \omterm{def:b3:complex-mechanisms:labile}{labil} atau inert dalam skala waktu satu hari?
\end{enumerate}

\medskip\noindent\textbf{Bagian II --- Sintesis.}
\begin{enumerate}[resume, beginpenalty=10000]
  \item Apakah yang dihasilkan $\ce{K2[PtCl4]}$ dengan dua \ce{NH3}? Mengapa?
  \item Rute langsung itu memberikan produk yang tidak murni. Rute Dhara mula-mula mengubah
        $\ce{[PtCl4]^{2-}}$ menjadi $\ce{[PtI4]^{2-}}$ dengan KI berlebih. Mengapa iodida berguna?
  \item Penambahan \ce{NH3} memberikan $cis$-$\ce{[PtI2(NH3)2]}$. Jelaskan stereokimianya.
  \item Iodida-iodida itu kemudian disingkirkan dengan perak nitrat dalam air. Apakah yang terbentuk?
  \item Penambahan KCl memberikan cisplatin. Mengapa konfigurasinya tidak berubah?
  \item Bagaimanakah transplatin dapat dibuat sebagai gantinya?
  \item Mengapa transplatin tidak berguna sebagai obat, meskipun zat itu juga bereaksi?
\end{enumerate}

\medskip\noindent\textbf{Bagian III --- Akuasi.}
\begin{enumerate}[resume, beginpenalty=10000]
  \item Hitunglah waktu paruh akuasi pertama.
  \item Hitunglah fraksi yang terakuasi setelah \qty{2.0}{h} dalam larutan bebas klorida.
  \item Mengapa kompleks akua merupakan bentuk yang reaktif terhadap DNA?
  \item Apakah akuasi itu disosiatif atau asosiatif? Bagaimanakah \omterm{def:b3:complex-mechanisms:activation-volume}{volume aktivasinya}?
  \item Mengapa larutan obat untuk suntikan dibuat dalam larutan garam fisiologis?
\end{enumerate}

\medskip\noindent\textbf{Bagian IV --- Plasma dan sel.}
\begin{enumerate}[resume, beginpenalty=10000]
  \item Tuliskan fraksi kesetimbangan kompleks akua sebagai fungsi
        $[\ce{Cl-}]$.
  \item Hitunglah fraksi itu dalam plasma.
  \item Hitunglah fraksi itu di dalam sel.
  \item Hitunglah rasionya.
  \item Mengapa, karena itu, obat tersebut terutama bekerja di dalam sel?
  \item Ligan-ligan lain manakah di dalam sel yang dapat mengikat platina dan menonaktifkan
        obat itu?
  \item Mengapa kedua ligan amin tetap terikat sepanjang waktu?
  \item Nyatakan hasilnya: rasio fraksi kompleks akua di dalam sel terhadap
        fraksinya dalam plasma.
\end{enumerate}
\end{problem}
